To complete the proof, let $r: \mathbb{V} \longrightarrow AT$ be a function. By the axiom of choice (but see exercise 1) there exists a function $g: \mathbb{V} \longrightarrow A\Omega$ such that $g.A\rho = r$. Since $g^{\#}.A\rho$ is an $\Omega$-homomorphism (by 2.3 and the first part of the proof), and the restriction of $g^{\#}.A\rho$ to $\mathbb{V}$ coincides with $r$, we have from 1.20 that $g^{\#}.A\rho = r^{\#}$. Since $g^{\#}$ maps equations into $E_A$ (2.4) and $A\rho$ identifies elements of $E_A$, $AT$ satisfies $E$ as desired.

$AT$ is called the free $(\Omega, E)$-algebra generated by $A$. The map $A\rho: A\Omega \longrightarrow AT$ presents $AT$ by “generators and relations,” and “free” means that there are just enough relations to satisfy $E$, but no more. To a category theorist, this intuitively correct formulation would be justified by the following result:

2.5 The Universal Property of $AT$. For each set $A$ define the insertion-of-the-variables map $A\eta: A \longrightarrow AT$ by $\langle a, A\eta \rangle = [a]$. Then for every $(\Omega, E)$-algebra $(X, \delta)$ and every function $f: A \longrightarrow X$ there exists a unique $\Omega$-homomorphism $f^{\#}: AT \longrightarrow (X, \delta)$ extending $f$.